% Original analytic teaching diagram, 2026-11-12.
% 1000 items from lines 甲, 乙, 丙 (500, 300, 200) with 10, 12, 20 defectives.
% (a) area diagram to scale: column width = share of line, dark strip height = defect rate
%     (2%, 4%, 10%); area of a strip = P(line and defective).
% (b) the 42 defectives alone, 1 mm per item: posterior shares 10/42, 12/42, 20/42.
\definecolor{VB}{HTML}{7188A3}\definecolor{VBF}{HTML}{E5EAF0}
\definecolor{VG}{HTML}{829C8B}\definecolor{VGF}{HTML}{E8EEE9}
\definecolor{VY}{HTML}{B4A16F}\definecolor{VYF}{HTML}{F1ECDD}
\definecolor{VR}{HTML}{A7787D}\definecolor{VRF}{HTML}{F0E5E6}
\definecolor{VP}{HTML}{9588A5}\definecolor{VPF}{HTML}{ECE8F0}
\definecolor{VI}{HTML}{3B4252}\definecolor{VM}{HTML}{6B7280}\definecolor{VL}{HTML}{4C4D4F}
\begin{tikzpicture}[x=1mm,y=1mm,>=Stealth,font=\figurefont\fontsize{8}{10}\selectfont,
lab/.style={text=FigureMuted,font=\figurefont\fontsize{7.5}{9}\selectfont},
txt/.style={text=VI},
cell/.style={draw=VL,line width=.6pt},
def/.style={draw=VB,fill=VB,line width=.5pt}]
\def\H{40}
% (a) area diagram, width 60 mm = 1000 items
\begin{scope}
\draw[cell] (0,0) rectangle (30,\H);
\draw[cell] (30,0) rectangle (48,\H);
\draw[cell] (48,0) rectangle (60,\H);
\fill[def] (0,0) rectangle (30,0.8);
\fill[def] (30,0) rectangle (48,1.6);
\fill[def] (48,0) rectangle (60,4.0);
\node[txt] at (15,\H+3) {甲};
\node[txt] at (39,\H+3) {乙};
\node[txt] at (54,\H+3) {丙};
\node[lab] at (15,\H/2) {500件};
\node[lab] at (39,\H/2) {300件};
\node[lab] at (54,\H/2) {200件};
\node[lab] at (15,-3) {10};
\node[lab] at (39,-3) {12};
\node[lab] at (54,-3) {20};
\node[txt] at (30,-9.5) {(a)};
\end{scope}
% (b) the defectives only, 1 mm per item
\begin{scope}[shift={(68,14)}]
\draw[def,fill=VBF] (0,0) rectangle (10,9);
\draw[def,fill=VBF] (10,0) rectangle (22,9);
\draw[def,fill=VBF] (22,0) rectangle (42,9);
\node[txt] at (5,4.5) {甲}; \node[txt] at (16,4.5) {乙}; \node[txt] at (32,4.5) {丙};
\node[txt] at (5,-3.5) {$\tfrac{10}{42}$};
\node[txt] at (16,-3.5) {$\tfrac{12}{42}$};
\node[txt] at (32,-3.5) {$\tfrac{20}{42}$};
\node[lab] at (21,12.5) {42件次品};
\node[txt] at (21,-23.5) {(b)};
\end{scope}
\end{tikzpicture}
